\section{Video Demo}

We demonstrate \dslr{} ability to generate precise static reconstructions and segment movable and moving objects through video demos. In the video demo at \cite{url1}, we visualize the results of the static reconstructions of MOVES and compare it with the best baseline. In this video demo at \cite{url2}, we compare the performance of our models on a the sparse industrial dataset - ARD-16. Our model is able to segment out the movable and moving objects effectively when compared to the best baseline. In this video demo at \cite{url3}, we visualize the segmentation of the movable and the moving objects in a LiDAR sequence.

 
\section{Conclusion}
In this work, we address the challenge that Autonomous Vehicles(AVs) face in the absence of labelled segmentation information. We try to close the gap between settings that have several labelled datasets available and several label-devoid settings. Non-stationary object segmentation and accurate static structures reconstruction are major research challenges in these settings. We present a novel GAN-based adversarial model, \dslr{} that is augmented with a robust pair discriminator. We utilize a smartly selected LiDAR triplet using hard-negative mining that serves two purposes - generate latent pairs for the discriminator training (and) augment the discriminator with contrastive loss to increase its discriminative capacity. Ablation studies show the benefit of this setting over a vanilla discriminator setting. \dslr{} segments out accurate \textit{movable} and moving objects using static reconstructions without using labelled information. The generated static structures that replace dynamic objects and corresponding occlusions supplement the available static structures which improves downstream task performance. Our work is of practical benefit for several industrial settings - that may not have labelled segmentation and object information due to several constraints. We also demonstrate that our method is adaptable to datasets that lack static-dynamic correspondence information and performs well in the absence of segmentation-labels. \dslr{} currently is not able to surpass the performance in the labelled settings - where labelleld segmentation information is available for navigation. We do perform better for long sequences in the presence of high dynamism and multiple loop closures. These are realistic and frequent scenarios that AVs come across. For moderate to low dynamism sequences, segmentation based methods perform better. A promising future work in this direction is to find more accurate static reconstruction methods to address this gap and perform better than baselines that use segmentation information. Another future work of immediate interest that we have planned to work on is to develop methods that can generate  paired correspondence information from dataset that do not have such correspondence pairs available (e.g. KITTI). This would enable us to eliminate the need for Domain Adaptation and boost static performance for existing unpaired datasets.

\section{Acknowledgment}
This work was supported by the Indian Institute of Technology, Delhi, India.